% Lösungen Einheit 4 von 5 – Grundwert berechnen (zusammenbau v0.1)
\einheitenkopf{Einheit 4 von 5 · Grundwert berechnen}

\erg{20}{a) Grundwert (gegeben: Prozentsatz $30\,\%$, Prozentwert $9$ Kinder)}
\erg{21}{a) $20$ Kästchen \quad b) $2 \cdot 35 = 70$ kg \quad c) $1\,\% = 4$ kg; $100\,\% = 400$ kg \quad d) $81 : 36 = 2{,}25$; $2{,}25 \cdot 100 = 225$ kg \quad e) $1\,\% = 0{,}28$; $100\,\% = 28$ Kinder \quad f) Prozentwert gesucht: $0{,}4 \cdot 65 = 26$ Kinder \quad g) $4 \cdot 4 = 16\,€$ (nicht $25\,\%$ von $4\,€$)}
\erg{22}{a) Fehler: Sven rechnet $30\,\%$ von $12$ kg, gesucht ist das Ganze. Richtig: $1\,\% = 0{,}4$ kg; $100\,\% = 40$ kg. \quad b) Nach dem Teilen hat man den Wert für $1\,\%$; das Ganze sind $100\,\%$, also hundertmal so viel. \quad c) $10\,\% = 4{,}5$ kg; Ganzes $10 \cdot 4{,}5 = 45$ kg \quad d) $1\,\% = 0{,}2$ h; $100\,\% = 20$ Stunden – ja, das reicht für $18$ Stunden.}
